\subsection*{一步一步建立直观}

\paragraph{1. 每一幅图都是一次计算的状态}

左图给出原始向量，中图是第一次动作的输出，右图才是最终结果。先让旋转角为零，只保留剪切，确认三幅图的含义。矩阵乘法不是把两个数表逐项相乘，而是把前一步产生的向量交给后一步处理。



\paragraph{2. 把中间状态代入第二条规则}

令 $B$ 是旋转，$A$ 是剪切。先算 $Bx$，再将这个新向量当成 $A$ 的输入，得到 $A(Bx)$。对每一根标准基箭头重复这个过程，就得到了乘积矩阵的各列。因此列的几何解释直接导出行列求和的计算公式。



\paragraph{3. 交换顺序是一项可检验的主张}

切换为先剪切再旋转，中图首先发生改变，最后的红色向量通常也不同。原因是第二次操作面对的方向已由第一次改变。结合律只改变分组，动作队列不变；交换律却重排了队列，两者不能混为一谈。



\subsection*{定义、命题与证明}

\paragraph{定义：映射的复合次序}

若 $B$ 把输入送到中间空间，$A$ 再作用于中间结果，则复合映射写成 $A\circ B$，矩阵写成 $AB$。最右侧的矩阵先作用。

\paragraph{命题：矩阵乘法的列规则}

复合矩阵 $AB$ 的第 $j$ 列是 $A$ 作用在 $B$ 的第 $j$ 列所得的向量。各矩阵的输入输出维度须匹配。

\begin{proof}

把标准基向量 $e_j$ 送入复合映射。第一步得到 $Be_j$，它就是 $B$ 的第 $j$ 列；第二步得到 $A(Be_j)$。

一个矩阵的第 $j$ 列也等于它作用在 $e_j$ 上的结果，所以复合矩阵的各列只能按此规则组成。交换 $A$、$B$ 会改变动作顺序，通常也改变结果。

\end{proof}
